% \dayonesegone/notation.tex
\begin{frame}
  \onslide<1->{\frametitle{Understanding Kalman Filter Notation}
\framesubtitle{Drilldown on State Dynamics}
}
 \begin{itemize}
\onslide<2->{ \item State Dynamics Model: $x(k)=\Phi x(k-1) + w(k)$}
\onslide<3->{ \item $x(k)$ is the \textit{\textbf{state vector}} with $n$
 elements, written as a column vector:}
\onslide<4->{ \item $x(k)=\state (k)=\statek$}
\end{itemize}
\onslide<5->{\center --------------------------------------------------------------------}
\end{frame}

\begin{frame}
  \onslide<1->{
  \frametitle{Understanding Kalman Filter Notation}
  \framesubtitle{Drilldown State Dynamics}
}
 \begin{itemize}
\onslide<2->{ \item Discrete State Dynamics}
\onslide<3->{ \item KF doesn't work for nonlinear dynamics function $f[x(k-1)]$ }
\onslide<4->{ \item Restriction to matrices:  $\simplestatedyneq$}
\onslide<5->{ \item $\stm{k}$ is called the \textbf{state transition matrix}.  It is 
$n\times n$ in size.}
\end{itemize}
\onslide<6->{\center --------------------------------------------------------------------}
\end{frame}

\begin{frame}
  \onslide<1->{
    \frametitle{Understanding Kalman Filter Notation}
    \framesubtitle{Drilldown on State Dynamics}
}
 \begin{itemize}
\onslide<2->{ \item State Transition Matrix: $\stmeqtwo{k}{k-1}$}
\onslide<3->{ \item Extrapolates state vector from time step $k-1$ to time step
 $k$ (pulling from past to present)}
\onslide<4->{ \item Indicates how state evolves over time without forcing input}
\onslide<5->{ \item Remember:  This is just a model!  Don't know
 true state, because noise forcing function generates disturbances!}
\end{itemize}
\onslide<6->{\center --------------------------------------------------------------------}
\end{frame}

\begin{frame}
  \onslide<1->{
    \frametitle{Understanding Kalman Filter Notation}
    \framesubtitle{Drilldown on State Dynamics}
}
 \begin{itemize}
\onslide<2->{ \item We could throw in other factors and terms}
 \begin{itemize}
\onslide<3->{ \item A matrix $\sdm{k}$ to distribute the Gaussian white noise
 to the states}
\onslide<4->{ \item An additional, deterministic control $u(k)$, and
 associated distribution matrix $\Lambda(k)$}
 \end{itemize}
\onslide<5->{ \item For simplicity, ignore control}
\onslide<6->{ \item Simplifying / decluttering notation
 possible}
\end{itemize}
\onslide<7->{\center --------------------------------------------------------------------}
\end{frame}

\begin{frame}
  \onslide<1->{
    \frametitle{The Kalman Filter Solution:}  
    \framesubtitle{State Estimate Extrapolation}
}
\begin{itemize}
\onslide<2->{\item Given an estimate $\stateest{k-1}{+}$ at time $k-1$, we extrapolate it forward in time 
by the state transition matrix:
\begin{equation}\label{stateextrapeq}
\stateextrapeq 
\end{equation}
}
\onslide<3->{\item Q: Why is white noise $w$ \textbf{not included} in extrapolation?}\\  
\onslide<4->{A:  Because white noise forcing function is \textbf{zero-mean}  
Difference equation is \textbf{homogeneous on average}, therefore we use homogeneous solution, i.e., 
the state transition matrix.}
\end{itemize} 
\onslide<5->{\center --------------------------------------------------------------------}
\end{frame}

\begin{frame}
  \onslide<1->{
    \frametitle{The Kalman Filter Solution:}
	\framesubtitle{Covariance Extrapolation}
}
\begin{itemize}
\onslide<2->{\item Q: How accurate is the solution?}\\
\onslide<3->{A: Answer given by \textbf{estimation error covariance matrix} $P(k)$ }
\onslide<4->{\item Like state estimate, covariance is extrapolated forward in time and then updated:}
\onslide<5->
{\begin{equation}\label{altcovextrapeq}
\altcovextrapeq
\end{equation}
where $Q(k)$ is the \textbf{process noise covariance matrix} }
\onslide<6->{\item $Q(k)$ measures dispersion in white noise forcing function $w(k)$}
\onslide<7->{\item $P^{-}(k)$ defined as expected estimation error: $\esterrdefminus$}
\end{itemize} 
\onslide<8->{\center --------------------------------------------------------------------}
\end{frame}

\begin{frame}
  \onslide<1->{
    \frametitle{The Kalman Filter Solution:}
	\framesubtitle{Covariance Extrapolation}
}
\begin{itemize}
\onslide<2->{\item What does Eq. \eqref{altcovextrapeq} really mean?}
\onslide<3->{\item $\stm{k}\covmat{k-1}{+}\stm{k}\trans$ represents effect of state dynamics 
on estimation error covariance}
\begin{itemize}
\onslide<4->{\item $P(k)$ is positive definite and $\Phi$'s action on $P(k)$ preserves positive definiteness }
\end{itemize}
\onslide<5->{\item $Q(k)$ increases uncertainty due to additive process noise}
\end{itemize} 
\onslide<6->{\center --------------------------------------------------------------------}
\end{frame}

